\subsection{COMPOSITION OF TWO CORRESPONDENCES}

Let \(G, G'\) be two graphs. Let A denote the set \(pr_1G\) and let C denote the set \(pr_2G'\) . The relation \((\exists y)((x, y) \in G\) and \((y, z) \in G')\) implies that \((x, z) \in A \times C\) , and therefore has a graph with respect to x and z.

DEFINITION 6. Let G, G' be two graphs. The graph (with respect to x and z) of the relation \((\exists y)((x, y) \in G \text{ and } (y, z) \in G')\) is called the composition of G' and G, and is denoted by G' o G (or sometimes by G'G).

PROPOSITION 3. Let G, G' be two graphs. The inverse of G' o G is then \(\overline{G} \circ \overline{G'}\) .

For the relation `` \((x, y) \in G\) and \((y, z) \in G'\) '' is equivalent to

\[
(z, y) \in \overline {{\mathbf {G} ^ {\prime}}} \text {   and   } (y, x) \in \overline {{\mathbf {G}}}
\]

PROPOSITION 4. Let \(G_{1}\) , \(G_{2}\) , \(G_{3}\) be graphs. Then

\[
\left(\mathrm{G} _ {3} \circ \mathrm{G} _ {2}\right) \circ \mathrm{G} _ {1} = \mathrm{G} _ {3} \circ \left(\mathrm{G} _ {2} \circ \mathrm{G} _ {1}\right).
\]

The relation \((x, t) \in (\mathbf{G}_3 \circ \mathbf{G}_2) \circ \mathbf{G}_1\) is equivalent to the relation

\[
(\exists y) ((x, y) \in \mathrm{G} _ {1} \text {   and   } \exists z ((y, z) \in \mathrm{G} _ {2} \text {   and   } (z, t) \in \mathrm{G} _ {3}))
\]

and therefore (by C33 (Chapter I, §4, no. 3)) to the relation

\[
(\exists y) (\exists z) ((x, y) \in \mathrm{G} _ {1} \text {   and   } (y, z) \in \mathrm{G} _ {2} \text {   and   } (z, t) \in \mathrm{G} _ {3}).\tag{1}
\]

Similarly, the relation \((x, t) \in \mathbf{G}_3 \circ (\mathbf{G}_2 \circ \mathbf{G}_1)\) is equivalent to

\[
(\exists z) (\exists y) ((x, y) \in \mathrm{G} _ {1} \text {   and   } (y, z) \in \mathrm{G} _ {2} \text {   and   } (z, t) \in \mathrm{G} _ {3}).\tag{2}
\]

But the relations (1) and (2) are equivalent; hence the result.

¶ The graph \(G_{3} \circ (G_{2} \circ G_{1})\) is denoted by \(G_{3} \circ G_{2} \circ G_{1}\) . Similarly, if \(G_{1}, G_{2}, G_{3}, G_{4}\) are graphs, we put

\[
\mathrm{G} _ {4} \circ (\mathrm{G} _ {3} \circ \mathrm{G} _ {2} \circ \mathrm{G} _ {1}) = \mathrm{G} _ {4} \circ \mathrm{G} _ {3} \circ \mathrm{G} _ {2} \circ \mathrm{G} _ {1};
\]

and so on.

PROPOSITION 5. Let G, G' be graphs and let A be a set. Then

\[
(\mathbf {G} ^ {\prime} \circ \mathbf {G}) \langle \mathbf {A} \rangle = \mathbf {G} ^ {\prime} \langle \mathbf {G} \langle \mathbf {A} \rangle \rangle .
\]

For by virtue of C33 (Chapter I, § 4, no. 3) the relation \(z \in (\mathbf{G}' \circ \mathbf{G})\langle \mathbf{A} \rangle\) is equivalent to

\[
(\exists y) ((\exists x) (x \in A \text {   and   } (x, y) \in G) \text {   and   } (y, z) \in G ^ {\prime})
\]

and is therefore equivalent to \((\exists y)(y \in G \langle A \rangle\) and \((y, z) \in G')\) ; hence the result.

If G and \(G'\) are two graphs, we have \(\mathrm{pr}_{1}(\mathrm{G}' \circ \mathrm{G}) = \mathrm{G}\langle\mathrm{pr}_{1}\mathrm{G}'\rangle\) , and \(\mathrm{pr}_{2}(\mathrm{G}' \circ \mathrm{G}) = \mathrm{G}'\langle\mathrm{pr}_{2}\mathrm{G}\rangle\) . For example, to prove the second of these relations it is enough to note that the relation \(z \in \mathrm{pr}_2(G' \circ G)\) is equivalent to \((\exists x)((x, z) \in G' \circ G)\) and therefore to

\[
(\exists y) ((\exists x) ((x, y) \in G) \quad \text { and } \quad (y, z) \in G ^ {\prime});
\]

but this is equivalent to \(z \in G' \langle pr_{2}G \rangle\) .

If G is a graph and X a set such that \(X \subset pr_{1}G\) , we have

\[
\mathbf {X} \subset \overline {{\mathbf {G}}} \langle \mathbf {G} \langle \mathbf {X} \rangle \rangle .
\]

For the relation \(x \in \mathbf{X}\) implies by hypothesis that \((\exists y)((x, y) \in \mathbf{G})\) ; but \((x, y) \in \mathbf{G}\) is equivalent to \((y, x) \in \overline{\mathbf{G}}\) , and on the other hand \((x, y) \in \mathbf{G}\) implies \((\exists z)(z \in \mathbf{X} \text{ and } (z, y) \in \mathbf{G})\) ; hence \(x \in \mathbf{X}\) implies

\[
(\exists y) ((\exists z) (z \in X \text { and } (z, y) \in G) \text { and } (y, x) \in \overline {{G}} ^ {- 1}),
\]

that is to say, \(x \in \overline{\mathbf{G}}\langle \mathbf{G}\langle \mathbf{X}\rangle\rangle\) .

It is clear that if \(G_{1}\) , \(G_{2}\) , \(G_{1}^{\prime}\) , \(G_{2}^{\prime}\) are graphs, the relations \(G_{1} \subset G_{2}\) and \(G_{1}^{\prime} \subset G_{2}^{\prime}\) imply \(G_{1}^{\prime} \circ G_{1} \subset G_{2}^{\prime} \circ G_{2}\) .

Let \(\Gamma = (G, A, B)\) and \(\Gamma' = (G', B, C)\) be two correspondences such that the target of \(\Gamma\) is the same as the source of \(\Gamma'\) . From the above discussion we have \(\mathrm{pr}_1(G' \circ G) \subset \mathrm{pr}_1G \subset A\) and \(\mathrm{pr}_2(G' \circ G) \subset \mathrm{pr}_2G' \subset C\) ; hence we may state the following definition:

DEFINITION 7. Let \(\Gamma = (G, A, B)\) and \(\Gamma' = (G', B, C)\) be two correspondences such that the target of \(\Gamma\) is the source of \(\Gamma'\) . Then the correspondence \((G' \circ G, A, C)\) is called the composition of \(\Gamma'\) and \(\Gamma\) , and is denoted by \(\Gamma' \circ \Gamma\) (or sometimes \(\Gamma' \Gamma\) ).

It follows immediately from Proposition 5 that if X is a subset of A we have \((\Gamma'\circ\Gamma)\langle X\rangle=\Gamma'\langle\Gamma\langle X\rangle\rangle\) . Furthermore, since the target of \(\overline{\Gamma}^{1}\) is the same as the source of \(\overline{\Gamma}\) , the inverse of \(\Gamma'\circ\Gamma\) is \(\overline{\Gamma}\circ\overline{\Gamma}'\) , by Proposition 3.

DEFINITION 8. If A is a set, the set \(\Delta_{\mathbf{A}}\) of all objects of the form \((x, x)\) , where \(x \in \mathbf{A}\) , is called the diagonal of \(\mathbf{A} \times \mathbf{A}\) .

Clearly we have \(\mathrm{pr}_1\Delta_{\mathbf{A}} = \mathrm{pr}_2\Delta_{\mathbf{A}} = \mathbf{A}\) . The correspondence

\[
\mathbf {I} _ {\mathbf {A}} = (\Delta_ {\mathbf {A}}, \mathbf {A}, \mathbf {A})
\]

is called the identity correspondence on A; it is its own inverse.

If \(\Gamma\) is a correspondence between A and B and if \(I_{A}, I_{B}\) are the identity correspondences on A, B, respectively, then \(\Gamma \circ I_{A} = I_{B} \circ \Gamma = \Gamma\) .

